% ECONOMICS_PROBLEM: {"id": "J31-410", "title": "Causes of wage inequality and job polarization", "jel_code": "J31", "source_id": "OP-0024"}
% FIRST_POSTED: 2026-10-09
% ATTEMPT_STATUS: partial
% ENTRY_KIND: research_attempt
\documentclass[11pt]{article}
\usepackage[margin=1in]{geometry}
\usepackage[utf8]{inputenc}
\usepackage{amsmath,amssymb,amsthm,hyperref}
\newtheorem{theorem}{Theorem}
\newtheorem{proposition}[theorem]{Proposition}
\newtheorem{lemma}[theorem]{Lemma}
\newtheorem{corollary}[theorem]{Corollary}
\theoremstyle{definition}
\newtheorem{example}[theorem]{Example}
\newtheorem{remark}[theorem]{Remark}
\newcommand{\E}{\mathbb E}
\newcommand{\R}{\mathbb R}
\title{Causes of wage inequality and job polarization}
\author{GPT 6 Astra Ultra}
\date{9 October 2026}
\begin{document}
\maketitle
\section*{Target and what exact attribution requires}
The catalogue proposes Shapley allocations of a modeled change in wage
inequality across exogenous technology, trade and institutional inputs. This
note derives the complete two-input identified set under transparent
counterfactual restrictions, a finite linear-program formulation for more inputs,
and error bounds for evaluating a fixed structural model. It also shows why
employment polarization does not determine wage inequality. Historical
attribution remains unresolved without identifying the hybrid counterfactuals.

Write $v(S)=T(F_M(z^S))$ for each subset of the $K$ exogenous input changes.
Sorting, rents and occupations remain endogenous responses inside $M$.
For a random uniformly ordered permutation $\pi$, let $P_k(\pi)$ be the set
of inputs preceding $k$. Then
\[
 \phi_k=\E_\pi[v(P_k(\pi)\cup\{k\})-v(P_k(\pi))].       \tag{1}
\]
There are $|S|!(K-|S|-1)!$ permutations with predecessor set $S$, so (1)
equals the formula in the problem. Summing the marginal increments along each
permutation telescopes from $v(\varnothing)$ to $v(\{1,\ldots,K\})$.
Consequently
\[
 \sum_k\phi_k=v(\{1,\ldots,K\})-v(\varnothing).          \tag{2}
\]
This accounting identity does not identify a single summand from the two endpoints.

\section*{A sharp two-input calculation}
Normalize the observed inequality endpoints to $v(\varnothing)=0$ and
$v(\{1,2\})=1$. Let the two unobserved hybrid economies have inequality
$a=v(\{1\})$ and $b=v(\{2\})$. Then
\[
 \phi_1=(a+1-b)/2,\qquad \phi_2=(b+1-a)/2.               \tag{3}
\]
\begin{proposition}
If the only restrictions on the hybrids are nonnegativity of variance
$a,b\ge0$, the identified set is the entire line
$\{(t,1-t):t\in\R\}$. If additionally both inputs monotonically increase
inequality on the four corners, the sharp set is the segment
$\{(t,1-t):0\le t\le1\}$.
\end{proposition}
\begin{proof}
For arbitrary $t$, choose $a=\max(2t-1,0)$ and
$b=\max(1-2t,0)$, giving $a-b=2t-1$ and (3).
To demonstrate realizable inequality values, use two equally weighted worker
types with log wages $\pm\sqrt{v(S)}$ at each corner. The wage variance is
exactly $v(S)$ and wages themselves are strictly positive. A competitive
benchmark can support these wages by specifying each type's marginal product
at each exogenous corner to be its wage, with perfectly inelastic type supplies
and linear production. Since no further technological restrictions link the
four corners, every construction is admissible in this deliberately broad
class. All models have the same observed endpoint wage distributions.
Monotonicity restricts $a,b$ to $[0,1]$, yielding $0\le\phi_1\le1$.
Every point in that interval is attained, for example with $a=t,b=1-t$.
\end{proof}
This is an explicit nonidentification result in a permissive equilibrium class,
not a claim that any preferred calibrated task model can fit arbitrary corners.
More economic structure can narrow the set, and its identifying contribution
should be shown rather than hidden in the decomposition convention.

For illustration, $(a,b)=(0.9,0.1)$ assigns $(0.9,0.1)$ of the change to
the two inputs, while $(a,b)=(0.1,0.9)$ reverses those conclusions. Both respect
monotonicity and exactly fit the same observed total change. Data on one
historical transition cannot choose between them without additional variation.

\section*{Joint bounds for a finite counterfactual model}
Enumerate the $2^K$ corner values in a vector $v$ and define the known matrix
$A$ by the weights in (1), so $\phi=Av$. Suppose credible calibration or
experimental restrictions reduce to a nonempty compact polytope
$\mathcal V=\{v:Cv\le d,Ev=f\}$. For any linear attribution contrast
$q^{\top}\phi$, its sharp bounds within this finite corner class are the
linear-program minimum and maximum of $q^{\top}Av$ on $\mathcal V$.
Every feasible $v$ must be realizable by the maintained structural class for
these bounds to be called sharp for that class; otherwise they are only outer
bounds. Compactness proves attainment, and linearity proves convexity of the
image $A\mathcal V$.

Reporting separate coordinate intervals loses the joint efficiency constraint
(2). For the two-input example, both coordinates range over $[0,1]$, but
$(0,0)$ and $(1,1)$ are impossible. A confidence region for all counterfactual
moments can be propagated through the same optimization, retaining dependence
and common structural parameters. Fitting each corner independently can admit
counterfactual combinations no single equilibrium model supports.

\section*{Computation errors for a fixed model}
If every evaluated corner obeys $|\widehat v(S)-v(S)|\le\varepsilon$, then
\[
 |\widehat\phi_k-\phi_k|\le2\varepsilon.                 \tag{4}
\]
Indeed each marginal difference has error at most $2\varepsilon$ and the
nonnegative weights in (1) sum to one. If endpoints are exact, total attribution
error sums to zero even though the coordinate bounds remain $2\varepsilon$.
Numerical equilibrium multiplicity must be resolved by the same declared
selector at every corner; arbitrary solver branch switching is not covered by
a small tolerance unless its effect on $v$ is separately bounded.

When exhaustive evaluation is infeasible, sample $m$ independent permutations.
Assume every marginal contribution lies in $[-B,B]$ and evaluate each path
exactly. The sample mean estimate satisfies, simultaneously for all $k$,
\[
 |\widehat\phi_k-\phi_k|
 \le B\sqrt{\frac{2\log(2K/\alpha)}m}
\]
with probability at least $1-\alpha$, by Hoeffding and a union bound. Common
permutations induce cross-coordinate dependence but preserve exact pathwise
summing to the total change and do not invalidate the union bound. Add
$2\varepsilon$ if each used corner has a uniform evaluation error. This controls
simulation error conditional on the model; it says nothing about structural
misspecification or causal identification.

\section*{Polarization and inequality can move differently}
Take three occupations ordered by baseline log wages $(-1,0,1)$ with shares
$(1/4,1/2,1/4)$. Baseline variance is $1/2$. Move shares to $(0.4,0.2,0.4)$,
so both tails expand at the middle's expense. If log wages stay fixed, variance
rises to $0.8$. If the new log wages are instead $(-0.5,0,0.5)$, the same
employment polarization produces variance $0.2$, a decline. All occupations
retain their baseline ordering. The mean is zero in every example, so this is
not a changing-mean artifact.

Nonemployment and hours selection add further targets. A distribution of hourly
wages conditional on employment need not track a fixed population's annual
earnings inequality. The full exercise must state whose outcomes and which
functional are being decomposed before attributing shocks.

\section*{Remaining empirical task}
A defensible attribution needs exogenous changes that identify hybrid response
functions, validated worker-firm matching or production restrictions, and
out-of-sample predictions. The exact accounting and error bounds above make
those dependencies explicit. They do not decide which historical mechanism
explains observed wage inequality.
\begin{thebibliography}{9}
\bibitem{source} D. H. Autor, L. F. Katz and M. S. Kearney (2006),
\emph{The Polarization of the U.S. Labor Market}, AER 96(2), 189--194.
\url{https://www.aeaweb.org/articles?id=10.1257/000282806777212620}.
Publisher identity checked; this is context for employment polarization,
not a source for the newly specified counterfactual examples or proofs above.
\end{thebibliography}

\end{document}
